\section{Complete services and actual policy costs}\label{sec:study67}
We evaluate the original nonlinear investment economy, not a finite-state replacement. The new design adds stronger prediction controls, prospective reuse, direct continuous-law policy inference and the existing two-control extension. The earlier tensor-reference and query studies remain separate evidence: in particular, the twenty strictly worse pure ReLU policies and the conventional winners of the earlier cold-start catalogue are not removed or reclassified.

\subsection{Frozen design and corrected accounting}
Two scalar tasks have $(d,T,\varepsilon)=(2,2,1/128)$ and $(8,3,1/64)$. Each returned policy uses twenty-four-bit downward state acquisition, thirty-two-bit downward actions, the original law, and the same uniform original-optimum budget. The action-label catalogue has ninety-six states per seed. ReLU, full quadratic, linear, additive linear spline and nearest-action fits share those labels for each of eight fixed seeds. Adaptive parabolic search, bisection, the upper-chord minimizer, an analytic terminal-derivative proposal and a midpoint rule supply five unfitted controls. Every arm uses the same pre-query screen, protected refinement and terminal kernel.

Two process repetitions give 180 comparison services. Another 24 services train once and load the saved model for twenty-four new workloads. Two inference services compare four complete callable controllers on 8,192 common-path rows per task. Six vector services cover two state dimensions and three horizons. The resulting catalogue has 212 predeclared processes, with a 900-second deadline per process. A failed or timed-out process remains part of its catalogue and work account. Repeated clocks are not independent training seeds.

The present execution corrects an error found by an independent rational trajectory check. A point interval used to initialize accumulated cost shared its lower and upper arrays. Updating these arrays in place overwrote the running lower endpoint with an upper endpoint. The repair allocates distinct arrays; it changes neither the Bellman kernels nor the policy-selection rule. We preserved the original execution and froze a complete corrected rerun with its own process receipts. The same random indices isolate the repair, but do not create another independent statistical sample. The active tables use only the corrected execution. Supplement, Section~\ref{supp:accounting67}, gives the counterexample and the reproduction boundary.

\subsection{Stronger predictors and observed certificate progress}
\begin{table}[htbp]
\centering
\caption{Ten selectors at the same original-optimum contract}\label{tab:comparison67}
{\small\setlength{\tabcolsep}{4pt}
\begin{tabular}{rlrrrrr}
\toprule
$d/T$ & Selector & Runs & Queries & Fit (s) & Full (s) & Warnings \\
\midrule
2/2 & Adaptive & 2 & 1234 & 0.000 & 0.365 & 0 \\
2/2 & Bisection & 2 & 1324 & 0.000 & 0.365 & 0 \\
2/2 & Upper chord & 2 & 1274 & 0.000 & 0.365 & 0 \\
2/2 & ReLU & 16 & 1209 & 0.368 & 0.766 & 6 \\
2/2 & Quadratic & 16 & 1219 & 0.024 & 0.365 & 0 \\
2/2 & Linear & 16 & 1224 & 0.023 & 0.365 & 0 \\
2/2 & Spline & 16 & 1214 & 0.023 & 0.365 & 0 \\
2/2 & Nearest & 16 & 1244 & 0.023 & 0.365 & 0 \\
2/2 & Derivative & 2 & 1284 & 0.000 & 0.365 & 0 \\
2/2 & Midpoint & 2 & 1289 & 0.000 & 0.365 & 0 \\
8/3 & Adaptive & 2 & 14469 & 0.000 & 1.567 & 0 \\
8/3 & Bisection & 2 & 17578 & 0.000 & 1.944 & 0 \\
8/3 & Upper chord & 2 & 14913 & 0.000 & 1.643 & 0 \\
8/3 & ReLU & 16 & 14569 & 0.919 & 2.620 & 32 \\
8/3 & Quadratic & 16 & 14569 & 0.526 & 2.169 & 0 \\
8/3 & Linear & 16 & 14638 & 0.523 & 2.144 & 0 \\
8/3 & Spline & 16 & 14663 & 0.525 & 2.219 & 0 \\
8/3 & Nearest & 16 & 14809 & 0.527 & 2.194 & 0 \\
8/3 & Derivative & 2 & 15204 & 0.000 & 1.769 & 0 \\
8/3 & Midpoint & 2 & 16383 & 0.000 & 1.793 & 0 \\
\bottomrule
\end{tabular}}
\par\smallskip
{\footnotesize Targets are $1/128$ for $d/T=2/2$ and $1/64$ for $8/3$. Full time is the parent launch-through-join receipt. Fitted selectors use eight declared seeds and two process repetitions; unfitted controls have two repetitions. Fit includes labels. Query and time entries are medians, not optimizer-success probabilities. Warning counts retain every executed fit. All arms use the same outward pre-query screen.}
\end{table}

The comparison charges complete process work from launch through shutdown and join. This includes imports, action labels, fitting, predictions, all Bellman queries, interval path evaluation, serialization and durable output. The terminal derivative is an analytic proposal in the original model, not an oracle for the preceding dates. Nearest action stores the finite training sample, not a policy on a tensor state grid. The upper-chord minimizer is supplied as an explicit conventional competitor rather than reserved for the learned arm.

All fitting seeds and warnings are reported. The table does not infer a population optimizer-success probability from a finite seed catalogue or a statistically identified runtime ranking from closely spaced clock measurements. The supplementary tables give every seed, every routed and declined proposal, and the separate changes in the upper witness and global lower bound. A routed query may improve neither endpoint; those observations remain in the progress account. The exact identity
\[
 (U_{\rm old}-L_{\rm old})-(U_{\rm new}-L_{\rm new})
 =(U_{\rm old}-U_{\rm new})+(L_{\rm new}-L_{\rm old})
\]
organizes the diagnostics. The margins in Theorem~\ref{thm:work-loss67} describe the work relevant to certification; the fitted action-label mean squared error is reported separately and is not asserted to control this identity.

\subsection{Measured reuse, including a declined model}
\begin{table}[htbp]
\centering
\caption{Prospective twenty-four-block model reuse}\label{tab:reuse67}
{\small\setlength{\tabcolsep}{4pt}
\begin{tabular}{rlrrcrr}
\toprule
$d/T$ & Selector & Runs & Full (s) & Full range (s) & Wins & Declines \\
\midrule
2/2 & Adaptive & 2 & 2.269 & 2.269--2.269 & 0/2 & 2 \\
2/2 & Bisection & 2 & 2.420 & 2.370--2.469 & 0/2 & 2 \\
2/2 & ReLU & 4 & 2.620 & 2.620--2.720 & 0/4 & 4 \\
2/2 & Quadratic & 4 & 2.269 & 2.269--2.269 & 4/4 & 4 \\
8/3 & Adaptive & 2 & 19.688 & 19.664--19.711 & 0/2 & 2 \\
8/3 & Bisection & 2 & 22.766 & 22.665--22.866 & 0/2 & 2 \\
8/3 & ReLU & 4 & 22.068 & 21.467--22.621 & 0/4 & 4 \\
8/3 & Quadratic & 4 & 20.967 & 20.717--21.121 & 0/4 & 4 \\
\bottomrule
\end{tabular}}
\par\smallskip
{\footnotesize Full time includes imports, fitting, model serialization, all loads and identity checks, twenty-four new workloads, the stale-contract fallback, all durable records and process exit. Wins compare matched repetition numbers with adaptive search; trained methods retain both declared seeds. Each service declines block index 7. Prefix crossings are separately recorded and are not labeled full-process break-even times.}
\end{table}

The reuse catalogue fixes ReLU and quadratic seeds 6603 and 6607 before its workloads. Each service fits and serializes once, then loads and checks the model for every block. Block index 7 deliberately presents a stale acquisition-contract identifier. The model is declined and the same conventional verifier is used. The file is not secretly repaired or silently accepted. Model loading, hashing, rejection, fallback, all queries and durable records belong to the measured service.

The supplementary crossing table compares each process with the adaptive process bearing the same repetition number. It reports the first observed cumulative crossing, any crossing that persists through the remaining observed blocks, and the final complete-process comparison. Prefix clocks end at a block's durable record and exclude the final summary and process exit; they are not mislabeled as complete costs of unexecuted stopping services. No crossing beyond the twenty-four-block window is extrapolated. This design replaces an assumed equal-query amortization ratio by actual repeated use, including unfavorable repetitions.

\subsection{Expected costs of the new callable policies}
\begin{table}[htbp]
\centering
\caption{Expected cost of the actual callable controllers}\label{tab:cost67}
{\small\setlength{\tabcolsep}{4pt}
\begin{tabular}{rlc}
\toprule
$d/T$ & Controller & Expected-cost interval \\
\midrule
2/2 & Adaptive & [0.55849, 0.62153] \\
2/2 & Bisection & [0.55849, 0.62153] \\
2/2 & Quadratic & [0.55845, 0.62149] \\
2/2 & ReLU & [0.55845, 0.62149] \\
8/3 & Adaptive & [0.67503, 0.72574] \\
8/3 & Bisection & [0.67503, 0.72573] \\
8/3 & Quadratic & [0.67501, 0.72572] \\
8/3 & ReLU & [0.67501, 0.72572] \\
\bottomrule
\end{tabular}}
\par\smallskip
{\footnotesize There are 8,192 distinct IID-model common-path rows per task, with continuous 48-bit initial and innovation bins and full tail bounds on ambiguous acquisition. The simultaneous family error is $1/100$ across both tasks, all absolute costs and paired differences. Endpoints are rounded outward. Repeated R66/R67 streams are counted only once.}
\end{table}

The inference experiment conditions on the frozen adaptive, bisection, ReLU and quadratic controllers, with fitted seed 6603 fixed in advance. Initial states are uniform on the full state cube. Each uniform innovation retains its original continuous law. Forty-eight-bit bins represent intervals of true continuous states and shocks, not point-mass substitutes. If a propagated state interval straddles an acquisition boundary, the unresolved tail receives its full valid cost-support enclosure. Such a path is neither discarded nor evaluated at its midpoint for statistical inference.

Four absolute costs and six paired differences are evaluated at each task. The bounded sample-variance inequality allocates $1/4000$ to each of forty endpoint tails, giving family error $1/100$ across both tasks. The conservative logarithm bound is nine, verified by an exact rational partial sum of the exponential series. There are 16,384 distinct common-path rows, not that number multiplied by the four methods or by the corrected rerun. A fixed pseudorandom stream ensures reproduction; independence is the explicitly stated bin-sampling model.

All twelve paired difference intervals contain zero. The table therefore supplies direct costs for the new complete callable policies without claiming that a change in query count identifies an economic gain. The lower and upper sample endpoints, negative paired gains and acquisition ambiguities are retained. The full paired intervals are in Supplement, Table~\ref{tab:differences67}. The earlier cost comparisons concern different policies and are not substituted for these observations.

\subsection{Two controls, two innovations and horizon cost}
\begin{table}[htbp]
\centering
\caption{Coupled two-control, two-innovation query scaling}\label{tab:vector67}
{\small\setlength{\tabcolsep}{4pt}
\begin{tabular}{rrrrrrrr}
\toprule
$d$ & $T$ & Full (s) & Queries & Descendants & Batch & RSS (KiB) & Recovery \\
\midrule
2 & 2 & 0.315 & 44 & 20 & 4 & 39348 & 0 \\
2 & 3 & 7.380 & 2957 & 1224 & 16 & 39328 & 0 \\
2 & 4 & 530.165 & 8106242 & 1894012 & 1024 & 43688 & 0 \\
8 & 2 & 0.365 & 44 & 20 & 4 & 39336 & 0 \\
8 & 3 & 11.491 & 2957 & 1224 & 16 & 39528 & 0 \\
8 & 4 & 841.366 & 8106112 & 1893976 & 1024 & 44836 & 0 \\
\bottomrule
\end{tabular}}
\par\smallskip
{\footnotesize Each service evaluates four designed initial states (all coordinates 0, 1, $1/4$ or $3/4$), with continuous-bin innovations, at target $1/4$. No state lattice or trained two-output selector is loaded. Batch is the largest recorded query batch, not the total number of live descendant states. RSS is process peak memory. Horizon-four cost is retained rather than extrapolated from the shorter horizons.}
\end{table}

The vector services use the original state-dependent capacity, the nonseparable action interaction and two independent continuous innovations. Each evaluates four designed initial states whose coordinates are all zero, one, one quarter or three quarters. These are declared scaling paths, not a four-observation estimate of a population policy cost. A triangular action cover replaces the scalar action interval; no state lattice or stored reference policy is consulted.

The horizon-four rows make the descendant cost visible. Query batches and peak resident memory are distinct records: the algorithm can allocate descendants before batching their evaluation. The uniform policy guarantee follows from the theorem and its precision conditions, not from extrapolation from four paths. This experiment executes the multivariate-control verifier; it is not a trained two-output neural comparison or a demonstration that long-horizon work is negligible.

\subsection{The purchase decision}
For a candidate complete service $j$ and benchmark $b$, let $[l_j,u_j]$ be the simultaneous interval for $J^b-J^j$, let $\Delta W_j$ be its measured work difference, and let $f_j$ be the incremental adoption fee. If $v\ge0$ values one normalized unit of expected cost and $p\ge0$ prices one unit of measured work, the net-benefit interval is
\begin{equation}\label{eq:purchase67}
 [\,v l_j-p\Delta W_j-f_j,\ v u_j-p\Delta W_j-f_j\,].
\end{equation}
A positive lower endpoint licenses adoption; a negative upper endpoint rules it out; otherwise the decision remains unresolved on the stated information. Applying this criterion to complete reuse receipts is different from treating a prefix crossing as a completed economic service. For a purchaser requiring only the stated accuracy contract, the lowest measured complete work among eligible services is a separate finite-catalogue criterion. Both are expressed in the original normalized units. No unobserved installation price, dollar calibration or empirical welfare magnitude is introduced.
